% !TEX program = xelatex
\documentclass[aspectratio=169,17pt]{beamer}

\usepackage{fontspec}
\usepackage{xeCJK}
\usepackage{booktabs}
\usepackage{hyperref}

\setCJKmainfont[Path=/Users/eric/Library/Fonts/]{LXGWWenKai-Regular.ttf}
\setmainfont[Path=/System/Library/Fonts/]{Helvetica.ttc}

\usetheme{Madrid}
\usecolortheme{default}
\setbeamertemplate{navigation symbols}{}
\setbeamertemplate{footline}[frame number]

\title{Snake Terminal}
\subtitle{双人贪吃蛇 —— 课程答辩}
\author{编程实战作业}
\date{\today}

\begin{document}

\begin{frame}
  \titlepage
\end{frame}

\begin{frame}{目录}
  \tableofcontents
\end{frame}

\section{项目概述}

\begin{frame}{项目简介}
  \begin{itemize}
    \item \textbf{Snake Terminal}：C++17 + raylib 的双人图形化贪吃蛇
    \item 场景驱动主循环；玩法可配置；含动画、音效与日志
    \item 构建：CMake + \texttt{build.sh}/\texttt{build.bat}；配置：\texttt{snake.cfg}
    \item 依赖：raylib（子模块）、miniaudio、dr\_mp3
  \end{itemize}

  \vspace{0.8em}
  \begin{block}{目标}
    逻辑与表现解耦、场景可扩展，在双人玩法之上做出完整的工程闭环。
  \end{block}
\end{frame}

\section{系统架构}

\begin{frame}{模块划分}
  \begin{columns}[T]
    \begin{column}{0.52\textwidth}
      \begin{itemize}
        \item \texttt{scene/} 菜单、配置、对局、结算
        \item \texttt{game/} 蛇逻辑 + 动画状态机
        \item \texttt{event/} 可消费事件分发
        \item \texttt{render/} Sprite / 分层 DrawList
        \item \texttt{audio/} 预解码混音音频
        \item \texttt{config/}、\texttt{utils/} 配置与日志
      \end{itemize}
    \end{column}
    \begin{column}{0.44\textwidth}
      \textbf{场景流}
      \begin{center}
        菜单 \\[0.3em]
        $\swarrow$\quad$\searrow$ \\[0.3em]
        配置\quad 游戏 \\[0.3em]
        $\searrow$\quad$\swarrow$ \\[0.3em]
        结算 $\rightarrow$ 菜单
      \end{center}
    \end{column}
  \end{columns}
\end{frame}

\begin{frame}{主循环与事件}
  \textbf{SceneManager::run}（单线程）

  \begin{enumerate}
    \item 按键白名单 $\rightarrow$ \texttt{InputEvent}
    \item \texttt{handle\_event}：通用钩子 $\rightarrow$ 类型分发（可 \texttt{consume}）
    \item \texttt{update(dt)} / 场景切换
    \item \texttt{render()}（BeginDrawing / EndDrawing）
  \end{enumerate}

  \vspace{0.5em}
  \begin{alertblock}{单线程约束}
    raylib 的 OpenGL context 与 \texttt{thread\_local} 要求窗口、纹理、绘制同线程。
  \end{alertblock}
\end{frame}

\begin{frame}{逻辑与动画分离}
  \begin{columns}[T]
    \begin{column}{0.48\textwidth}
      \textbf{Snake}（纯逻辑）
      \begin{itemize}
        \item 移动、碰撞、得分
        \item 复活、尾部切除
        \item 虚影 / 可碰撞开关
      \end{itemize}
    \end{column}
    \begin{column}{0.48\textwidth}
      \textbf{SnakeBody} 状态机
      \begin{itemize}
        \item Move：移动 / 加速特效
        \item Die：逐节 X 标记
        \item Waiting：虚影或冻结
      \end{itemize}
    \end{column}
  \end{columns}

  \vspace{0.8em}
  由 \texttt{GameScene} 驱动：\\
  MOVE $\to$ DIE $\to$ WAITING $\to$ MOVE \\
  双玩家独立 tick 计时；Shift 按下时 tick 间隔减半。
\end{frame}

\section{玩法与工程}

\begin{frame}{双模式与操作}
  \begin{table}
    \centering
    \small
    \begin{tabular}{@{}lll@{}}
      \toprule
      模式 & 规则 & 死亡后 \\
      \midrule
      Deathmatch & 撞墙 / 己 / 对方即终局 & 动画后结算 \\
      Time Race & 限时计分，可复活 & 切尾 + 虚影/重部署 \\
      \bottomrule
    \end{tabular}
  \end{table}

  \vspace{0.6em}
  P1：WASD\quad P2：IJKL\quad Shift：加速 \\
  ESC 暂停；暂停后 T 回菜单 / Y 强制结束
\end{frame}

\begin{frame}{可配置规则}
  \texttt{snake.cfg} 持久化，缺项自动补全：
  \begin{itemize}
    \item 音量、加速、环面地图、是否穿过对方
    \item 速度倍率、难度递增、限时与复活代价
    \item \textbf{虚影复活}：半透明预览位置后再部署
    \item \textbf{死亡动画打断}：阈值后可按键跳过（$-1$ 关闭）
  \end{itemize}
\end{frame}

\begin{frame}{技术亮点}
  \begin{enumerate}
    \item \textbf{场景 + 事件}：生命周期清晰，事件可拦截
    \item \textbf{逻辑 $\perp$ 表现}：Snake 与动画状态机解耦
    \item \textbf{自研音频}：预解码 PCM + 回调混音；BGM 独占 / SFX 并发池
    \item \textbf{日志}：级别过滤、彩色输出、接管 raylib 日志
    \item \textbf{上游贡献}：修复 raylib \texttt{ImageFromImage} 边界问题（已合并）
  \end{enumerate}
\end{frame}

\section{总结}

\begin{frame}[fragile]{构建与运行}
\begin{verbatim}
./build.sh                 # 增量构建
./build.sh --build-raylib  # 含 raylib
./dist/snake --loglevel debug
\end{verbatim}

  \vspace{0.5em}
  Windows 可用 \texttt{build.bat}；启动时切到可执行文件目录，保证相对路径正确。
\end{frame}

\begin{frame}{总结}
  \begin{block}{已完成}
    双人双模式；场景/事件框架；可配置玩法；动画与虚影；音频与日志。
  \end{block}

  \vspace{0.5em}
  \begin{block}{可继续}
    地图障碍 / 关卡；网络对战；资源池与更细的渲染拆分。
  \end{block}

  \vspace{1em}
  \begin{center}
    {\Large 谢谢老师！欢迎提问}
  \end{center}
\end{frame}

\end{document}
